\section{Revisão da Literatura}

%% Conflito: preservar ou converter + MT
% CITAÇÃO1 e CITAÇÃO2
A expansão agropecuária envolve uma escolha econômica entre usos alternativos da terra, na qual a conversão da vegetação pode ampliar a renda privada associada à produção, ao mesmo tempo em que reduz estoques de capital natural e serviços ecossistêmicos que nem sempre são incorporados aos retornos percebidos pelos agentes. Seus efeitos podem, portanto, ultrapassar os limites da propriedade convertida, manifestando-se na redução da regulação hídrica, na perda de solo e carbono e na redução da biodiversidade, com custos que recaem sobre outros agentes e territórios. No Cerrado brasileiro, por exemplo, a conversão da vegetação para pastagem foi associada a escoamento da água da chuva (\textit{surface runoff}) quase vinte vezes maior, enquanto a conversão para cana elevou substancialmente a perda de solo \cite{Schwamback2024}; na fronteira Amazônia--Cerrado, a substituição da vegetação por soja esteve associada à redução dos estoques de carbono mesmo em camadas profundas do solo \cite{Popin2025}. Tais efeitos ajudam a explicar por que decisões individualmente rentáveis podem produzir custos ambientais e sociais não refletidos no cálculo privado, distinção amplamente reconhecida sobre capital natural e serviços ecossistêmicos \cite{Brady2015,Mann2012,Feng2022,Vackarova2023}.

Com efeito, o problema permanece ainda menos resolvido quando comparam-se benefícios e custos privados e públicos, pois os serviços ambientais diferem entre territórios, operam em escalas distintas e muitos não possuem preços diretamente observáveis \cite{Crossman2009,Gascoigne2011,Schneibel2016}. Assim, a oposição entre conservar e produzir pode refletir não apenas um conflito tecnológico, mas diferenças entre retornos apropriados privadamente e custos ambientais distribuídos socialmente. Em Mato Grosso, onde a produção se expande sobre biomas distintos, torna-se ainda menos plausível supor que decisões semelhantes de conversão produzam custos econômicos e ambientais equivalentes no território. Neste sentido, estudos recentes apontam que os resultados são desiguais entre os biomas. 

Para a amazônia brasileira, \citeonline[p.~279]{Mann2012} estimam diretamente o descompasso entre retornos privados e serviços ecossistêmicos e estimam que: ``\emph{Based on the balance between the benefits and costs of conversion, failure to value ecosystem services reduces total social welfare by 7.13 billion dollars annually relative to an optimum}''. Para o Cerrado, estudos documentam efeitos relevantes da conversão sobre processos biofísicos, incluindo aumento do escoamento superficial, perda de solo e redução de estoques de carbono \cite{Schwamback2024,Popin2025}, mas tais efeitos não são necessariamente avaliados em uma métrica econômica comparável àquela utilizada para a Amazônia. 

No caso do Pantanal, existem tentativas de mensurar economicamente os serviços ambientais, mas os próprios resultados evidenciam dificuldades de escala, cobertura e comparabilidade. \cite{SeidlMoraes2000}, ao recalcularem para a Nhecolândia \footnote{A Nhecolândia é uma sub-região pantaneira localizada no estado de Mato Grosso do Sul, compreendida entre os rios Taquari e Negro \cite{silva1998delimitacao}.} estimativas globais de serviços ecossistêmicos com dados locais, obtiveram valor por hectare aproximadamente metade daquele derivado dos parâmetros globais, diferença atribuída à heterogeneidade biofísica regional. Mais recentemente, \cite{FairbrassEtAl2025}, ao estruturarem indicadores de capital natural para o Pantanal, identificaram lacunas justamente na mensuração da condição e do valor econômico do bioma, das emissões associadas às atividades produtivas e de serviços como ciclagem de nutrientes e biodiversidade. 

Além disso, mesmo esforços de harmonização entre os biomas brasileiros precisam excluir componentes não comparáveis e produzir estimativas para preencher lacunas existentes \cite{CampoliStivali2026}. Desta feita, essa assimetria de evidências limita comparações entre biomas e dificulta estabelecer em que medida um mesmo incentivo econômico à conversão produz consequências equivalentes em territórios ambientalmente distintos.

% CITAÇÃO3 e CITAÇÃO4
A par disso, a literatura converge em tratar a conversão da vegetação nativa como resultado de escolhas econômicas condicionadas por retornos relativos, direitos de propriedade e capacidade de enforcement \cite{Pfaff1999,ProbstEtAl2020,AssuncaoGandourRocha2023}. Esse enquadramento não supõe, porém, efeitos uniformes: a mesma regra pode produzir respostas distintas conforme variam os custos de oportunidade, a segurança fundiária (\textit{land tenure}) e a efetividade local da fiscalização. Permanece, então, pouco resolvida a identificação conjunta desses mecanismos, em parte porque os estudos apontam medidas e contextos institucionais distintos \cite{Sousa2016,SticklerEtAl2013}. Em Mato Grosso, onde sanções, embargos e custos de conformidade ao Código Florestal exibem eficácia desigual na contenção do desmatamento, essa lacuna torna-se especialmente relevante.

%% Ponte para H1 (Retornos, instituições e conversão) + MT
% as condições citadas no parágrafo anterior obedecem a incentivos institucionais
% CITAÇÃO5 e CITAÇÃO6
A produção acadêmica reconhece de forma majoritária que a conversão da vegetação resulta da interação entre retornos econômicos do uso da terra e instituições que condicionam sua tradução em decisões privadas. Preços, rentabilidade e acesso a mercados elevam a pressão de conversão, enquanto monitoramento, sanções críveis e restrições legais podem contê-la. Como sintetizam \citeonline[p.~1]{HardingEtAl2021}, ``\textit{the effect of environmental regulation depends on the economic pressure to use natural resources}''. As evidências, contudo, mostram que essa interação não produz respostas uniformes, a depender dos biomas e territórios.

Na Amazônia brasileira, Harding et al. (2021) estimam que a inclusão de municípios em lista prioritária reduziu o desmatamento associado aos preços das commodities, mas encontram também substituição entre cultivos e deslocamento da pressão para terras não protegidas sob outros instrumentos. Assunção et al. (2023), por sua vez, identificam efeito relevante do monitoramento por satélite e da fiscalização sobre o desmatamento amazônico, enquanto análises espaciais mostram que acessibilidade, aptidão produtiva e condições biofísicas também modificam a forma pela qual incentivos econômicos se convertem em uso da terra \cite{ChomitzThomas2003}. 

As divergências, portanto, parecem decorrer menos de resultados incompatíveis do que de diferenças entre instrumentos, períodos, escalas territoriais e mecanismos identificados: parte dos estudos isola preços, outros fiscalização, direitos de propriedade ou condições espaciais, dificultando estimar conjuntamente quanto da resposta à rentabilidade depende de cada componente institucional. Essa dificuldade é especialmente relevante em Mato Grosso, onde Sousa (2016) encontra forte capacidade dissuasória dos embargos, mas baixa efetividade potencial das multas diante de sua reduzida arrecadação. Assim, permanece em aberto não se retornos ou instituições importam, mas como sua interação determina a intensidade da conversão em territórios sujeitos simultaneamente a forte pressão produtiva e diferentes condições de enforcement, e até que ponto essa interação produz efeitos distintos entre biomas.


% FIM